\documentclass[10pt]{article}

% Shared presentation layer for the standalone R0--R7 development documents.
% Method equations remain authoritative in the canonical idea sketch.
\usepackage[a4paper,margin=18mm,headheight=14pt]{geometry}
\usepackage{fontspec}
\usepackage{xeCJK}
\xeCJKsetup{CJKspace=true}
\setmainfont{DejaVu Serif}
\setsansfont{DejaVu Sans}
\setmonofont{DejaVu Sans Mono}
\setCJKmainfont{Noto Serif CJK KR}
\setCJKsansfont{Noto Sans CJK KR}
\setCJKmonofont{Noto Sans Mono CJK KR}

\usepackage{amsmath,amssymb}
\usepackage{booktabs,longtable,tabularx,array}
\usepackage{enumitem}
\usepackage{xcolor}
\usepackage{hyperref}
\usepackage{fancyhdr}
\usepackage{microtype}
\usepackage[most]{tcolorbox}

\definecolor{CoreBlue}{HTML}{1F4E79}
\definecolor{CoreLight}{HTML}{EAF3F8}
\definecolor{StopRed}{HTML}{9C0006}
\definecolor{StopLight}{HTML}{FCE8E6}
\definecolor{GateGreen}{HTML}{38761D}
\definecolor{GateLight}{HTML}{EAF4E2}
\definecolor{DecisionOrange}{HTML}{B45F06}
\definecolor{DecisionLight}{HTML}{FFF2CC}

\hypersetup{
  colorlinks=true,
  linkcolor=CoreBlue,
  urlcolor=CoreBlue
}

\setlength{\parindent}{0pt}
\setlength{\parskip}{0.3em}
\setlength{\emergencystretch}{3em}
\setlist[itemize]{leftmargin=1.45em,itemsep=0.1em,topsep=0.2em}
\setlist[enumerate]{leftmargin=1.75em,itemsep=0.14em,topsep=0.22em}
\renewcommand{\arraystretch}{1.15}
\newcolumntype{L}[1]{>{\raggedright\arraybackslash}p{#1}}
\newcommand{\code}[1]{{\small\nolinkurl{#1}}}
\newcommand{\must}[1]{\textcolor{CoreBlue}{\textbf{#1}}}
\newcommand{\haltitem}[1]{\textcolor{StopRed}{\textbf{#1}}}
\newcommand{\passitem}[1]{\textcolor{GateGreen}{\textbf{#1}}}
\newcommand{\decisionitem}[1]{\textcolor{DecisionOrange}{\textbf{#1}}}

\newenvironment{contractbox}[1]
  {\begin{tcolorbox}[colback=CoreLight,colframe=CoreBlue,title={#1},fonttitle=\bfseries,before upper=\raggedright]}
  {\end{tcolorbox}}
\newenvironment{decisionbox}[1]
  {\begin{tcolorbox}[colback=DecisionLight,colframe=DecisionOrange,title={#1},fonttitle=\bfseries,before upper=\raggedright]}
  {\end{tcolorbox}}
\newenvironment{stopbox}[1]
  {\begin{tcolorbox}[colback=StopLight,colframe=StopRed,title={#1},fonttitle=\bfseries,before upper=\raggedright]}
  {\end{tcolorbox}}

\renewcommand{\DevChecklistDate}{2026-09-11}

\hypersetup{
  pdftitle={R2 Single-Case Global Overfit --- Wind3DGS Development Specification},
  pdfauthor={Wind3DGS Project}
}

\pagestyle{fancy}
\fancyhf{}
\lhead{\small Wind3DGS Development --- R2}
\rhead{\small Single-Case Global Overfit}
\cfoot{\thepage}

\title{\textbf{R2 --- Single-Case Global Overfit}\\
R2a fixed-seconds feasibility와 R2b curriculum/baseline}
\author{Wind3DGS Project}
\date{2026-08-22}

\begin{document}
\maketitle

\begin{contractbox}{문서 권위와 두 단계 원칙}
이 문서는 canonical sketch의 single-case Global 학습을 구현 단위로 분리한 파생 문서다.
\code{eq:rd-gaussian-token}, \code{eq:rd-latent-anchor},
\code{eq:rd-response-package}, \code{eq:rd-area-mass},
\code{eq:rd-attachment-whitening}, \code{eq:rd-global-response},
\code{eq:rd-exact-zoh}, \code{eq:rd-wind-pullback},
\code{eq:rd-mean-transport}, \code{eq:rd-response-loss},
\code{eq:rd-total-loss}, \code{eq:rd-teacher-period-spectrum}을 텍스트 기준으로 참조한다.
먼저 T1 measure/token initialization을 마친 뒤 curriculum이나 baseline 없이
\textbf{R2a fixed-seconds 최소 overfit}을 통과하고,
그 뒤에만 \textbf{R2b curriculum/baseline}을 실행한다.
\end{contractbox}

\DevChecklistPage{R2}{R1 선행 조건 대기}{
\DevProgressRow{0}{\hyperref[sec:r2-chapter-1]{1. 목적과 소유권}}{대기}{R1 통과 후 T1/R2a/R2b 실행 범위와 owner 확인}
\DevProgressRow{0}{\hyperref[sec:r2-chapter-2]{2. 진입 입력}}{대기}{Accepted label/probe/oracle와 초기 상태 계약 확보}
\DevProgressRow{0}{\hyperref[sec:r2-chapter-3]{3. 이미 동결된 계약}}{대기}{Package·loss·시간·evaluator 계약을 구현에 대조}
\DevProgressRow{0}{\hyperref[sec:r2-chapter-4]{4. 구현 체크리스트}}{대기}{T1, fixed-seconds overfit, 이후 curriculum/baseline}
\DevProgressRow{0}{\hyperref[sec:r2-chapter-5]{5. Open Design Decisions}}{대기}{Rank, loss weight, horizon과 capacity budget 동결}
\DevProgressRow{0}{\hyperref[sec:r2-chapter-6]{6. 산출물}}{대기}{Checkpoint/package/rollout 및 비용 report 생성}
\DevProgressRow{0}{\hyperref[sec:r2-chapter-7]{7. 필수 fixture와 metric}}{대기}{운동·위상·장기 안정성 및 oracle gap 검증}
\DevProgressRow{0}{\hyperref[sec:r2-chapter-8]{8. 종료 조건과 실패 경로}}{대기}{R2a 및 R2b 종료·schedule/fallback 판정}
}{
\begin{contractbox}{진입 경계}기존15개 development window와 P3 작은 굽힘 sample76개는 개발 근거로 보존한다. 사용자 지시에 따라 물리 수식 구현·검증 전 추가 데이터 발행은 보류한다. R1 oracle/GS·시작 상태 계약까지 확보한 뒤 R2를 실행한다. 학습 성능의 완료 근거는 아직 없다.\end{contractbox}
}

\tableofcontents
\newpage

\section{목적과 소유권}
\label{sec:r2-chapter-1}

R2의 핵심 질문은 ``현재 setup network와 response package가 믿을 수 있는 한 개의 teacher fixture를
실제로 외울 수 있는가''다. Generalization, multi-resplat consistency와 Local 필요성은 여기서 주장하지 않는다.

\begin{longtable}{L{0.20\linewidth}L{0.34\linewidth}L{0.36\linewidth}}
\toprule
단계 & 소유하는 질문 & 소유하지 않는 질문\\
\midrule
\endfirsthead
\toprule
단계 & 소유하는 질문 & 소유하지 않는 질문\\
\midrule
\endhead
R2a & 한 object/GS/material의 impulse/step을 fixed physical seconds에서 overfit하고 긴 free decay까지 안정한가 & curriculum 이득, baseline 우위, held-out generalization\\
R2b & R2a를 유지한 채 optional period curriculum이 equal-compute 이득을 내는가; 최소 baseline과 failure owner가 무엇인가 & Gate B의 object-disjoint claim\\
공통 & Global package prediction, exact construction, deterministic rollout, loss/metric contract와 trace & Local, selector, appearance/perceptual loss, broad dataset\\
\bottomrule
\end{longtable}

R2a는 학습 가능성을 가장 작은 폐쇄계에서 확인한다. R2b는 방법상의 선택과 비교를 정리하지만,
single-case baseline 승리는 논문 claim의 증거가 아니다. Learned Global held-out claim은 R4/Gate B가 소유한다.

\section{진입 입력}
\label{sec:r2-chapter-2}

\subsection{R0/R1 필수 artifact}

\begin{itemize}
  \item R0의 frozen setup/frame/package/state/tensor/coordinate/valid-set/serialization contract.
  \item R0의 \code{MetricRegistry}, \code{ThresholdRegistry}와 failure denominator.
  \item R1의 converged \code{TeacherPackage}, accepted teacher mesh/timestep, \code{CanonicalProbeMap},
  common-valid mask, mapping-noise floor와 base transport fixture.
  \item 동일 object/GS의 corrected \code{OracleGlobalPackage}, oracle response error와 exact-ZOH replay.
  \item Teacher/GS traction identity, SI unit, rank/normal branch와 source/config/artifact hash.
  \item \code{TeacherSpectrumStatus}: converged, 그리고 stable peak available/unavailable reason.
\end{itemize}

R1 종료 조건 중 하나라도 미충족이면 R2 training command를 실행하지 않는다.
현재 생성된 Newton development window 15개는 \code{training_eligible=false}이며 R2 label이 아니다.
후속 P3의76개 patch sample은 지정한 작은 굽힘 wind/reset의 수렴·원본·batch 검증을 통과했다.
\code{sample_scope_eligible=true}이며 큰 변형 또는 R1 전체 채택은 아니다.
이를 R2에 사용하려면 해당 response package의 oracle/GS 연결과 명시적 시작 상태 계약을 먼저 확인한다.
실제 accepted R1 artifact/hash와 eligibility를 모두 확인하고, loader의 개발용 opt-in으로 이 경계를 우회하지 않는다.
사용자가 후속 구조 모델로 선택한 P3 유한 회전 shell의 수식·solver 검사는 R1에서 진행한다.
선행 작은 굽힘 sample이나 새 수식 검사를 R2 학습 시작 승인으로 해석하지 않는다.

\subsection{R2 최소 데이터 slice}

\begin{itemize}
  \item 한 source object, 한 accepted independent GS, 한 material preset, 한 attachment.
  \item Impulse와 step-on/off response. Aero-off displaced free decay를 별도 sequence로 둔다.
  \item Prescribed wind/domain-wide \(\kappa_{\mathrm{aero}}\), zero-activation \(\tau_{\mathrm{guard}}\)와
  fixed \(\Delta t\), common probe/time grid.
  % [논문 작성 참고 -- 수정 이유]
  % R2 overfit에서 rho/C_n/C_parallel을 sample별로 조정하지 않고, R1이 동결한 단일 effective scale을 그대로 소비한다.
  \item Train-only overfit slice이지만 sample ID, sequence, initial state와 horizon을 manifest에서 고정한다.
  \item R2a는 \(T_{\rm dom}\)을 사용하지 않는 fixed-seconds \(H_{\rm smoke}\),
  \(H_{\rm long}\)과 stress horizon을 먼저 사용한다. 값은 Open Design Decision이다.
\end{itemize}

\subsection{Model 진입 상태}

사용자가 원하는 장면 길이는10--20 s이며, 현재 R1의 첫10 s 시간 간격 탐색은 추천 후보를
확보하지 못했다. 이를 accepted R2 label이나 확정된 \(H_{\rm long}\) 값으로 사용하지 않는다.
짧은 smoke와 최종 장기 사용 범위는 구분하고, R1 입력 적격성과 아래 horizon 동결 절차를 먼저 닫는다.

R1의 rest-start 변화 바람/velocity-reset 25개 trace는 개발 검산을 통과했지만 속도 수렴과 굽힘 상태
coverage에 한계가 있다. 적절성·수렴·입력 범위가 승인된 뒤
R2 slice 편입 여부를 판단한다. 이를 근거로 기존 impulse/step 및 aero-off displaced decay를 삭제하거나
현재 학습 적격성이 false인 reset trace를 필수 학습 입력으로 간주하지 않는다.

초기 target model은 patch encoder, fixed/capped latent token, token global context,
token-to-Gaussian equivariant query readout, positive area head와 area-derived mass,
Global response field와 positive pole/damping head로 제한한다. Setup에서 package를 한 번 cache하고
frame 동안 network를 다시 호출하지 않는다.

\section{이미 동결된 계약}
\label{sec:r2-chapter-3}

\subsection{범위와 exact construction}

\begin{itemize}
  \item R2에는 Global package와 deterministic exact-ZOH evaluator만 있다. Local field/selector/state,
  image/perceptual loss, optional SH residual과 test fine-tuning은 없다.
  \item Area는 positive이고 V0 mass는 \(m_i=(M_{\mathrm{ref}}/A_{\mathrm{ref}})a_i\)로 파생한다
  (sketch \code{eq:rd-area-mass}). Valid-set normalization은 frozen R0 policy를 그대로 쓴다.
  \item Attachment zero와 mass whitening은 penalty가 아니라 package construction으로 보장한다
  (sketch \code{eq:rd-attachment-whitening}).
  \item Rank/condition 검사, positive angular frequency와 modal damping ratio도 validator/construction이 소유한다.
  \item Learned \(\Omega_G,Z_G\)는 attachment projection과 mass whitening을 마친 최종 generalized
  coordinate의 pole이다. Raw-field coordinate의 pole을 basis construction 뒤 그대로 연결하지 않는다.
  \item Global은 sketch \code{eq:rd-global-response}의 independent damped response이고,
  frame-start load의 single exact ZOH만 사용한다 (sketch \code{eq:rd-exact-zoh}).
  \item Runtime wind는 GS에서 직접 \(F_i=a_i\tau_i\), \(Q_G=B_G^\top F\)로 계산한다.
  Common-probe map은 loss/evaluation에만 쓴다.
\end{itemize}

\subsection{Response objective의 고정 부분과 미결정 부분}

Teacher mode column은 sign/order/subspace gauge가 비유일하므로 primary label로 사용하지 않는다.
Sketch \code{eq:rd-response-loss}의 displacement, velocity, FRF amplitude/phase와 rollout response family를
frozen common-valid probe에서 학습한다. Position보다 velocity를 primary로 둔다.

그러나 \textbf{합/평균, reference-scale normalization, sequence reduction, FRF phase mask,
각 \(\lambda\)와 numeric threshold는 아직 동결되지 않았다}. 이 항목은 Open Design Decisions에서 선택하고
R2a acceptance run 전에 version/hash해야 한다. 단위가 다른 raw loss를 normalization 없이 더하지 않는다.

R2에서 sketch \code{eq:rd-total-loss}를 사용할 때 Local/pair/switch 항은 활성화하지 않는다.
T1 initialization에서는 local area/measure supervision을 활성화하고 patch-to-token path를 실제로
학습한다. Total-sum sanity만 loss로 쓰지 않는다. Relation auxiliary의 사용 여부와 T2에서 measure loss를
freeze/continue하는 선택은 development config지만, T1 phase 자체를 생략하지 않는다.
Total area/mass, positivity, attachment와 whitening처럼 construction으로 만족하는 조건을 중복 penalty로
성능 증거로 포장하지 않는다.

\subsection{Causal rollout와 physical time}

\begin{itemize}
  \item Network는 setup에서 한 번 package를 출력한다. Recurrent teacher forcing, scheduled sampling,
  target state injection과 opaque hidden-state correction을 사용하지 않는다.
  \item 각 rollout은 같은 explicit initial physical state에서 시작하고 evaluator의 \(q,\dot q\)만 전개한다.
  비정지/window 시작에는 명시적 state 초기화 또는 앞선 wind 구간 재생과 초기 재구성 오차 검증이 필요하다.
  Teacher 형상·속도를 저장한 사실만으로 package 내부 상태가 초기화되었다고 보지 않는다.
  Reset 경계를 일반 연속 window로 숨기거나 매 frame 정답 상태를 주입하지 않는다.
  \item Horizon 비교는 frame count가 아니라 seconds다. \(\Delta t\)가 달라도 같은 physical time을 비교한다.
  \item R2a fixed-seconds overfit을 먼저 통과하지 못하면 period curriculum을 열지 않는다.
  \item 학습 horizon보다 긴 in-envelope rollout의 NaN/divergence와 declared damping/control decay-order 붕괴는
  짧은 결과로 가리지 않는다.
\end{itemize}

\subsection{R2b curriculum truth table}

R1의 teacher spectrum 자체가 수렴하지 않았다면 R2 진입 실패다. Spectrum이 수렴한 뒤에는
sketch \code{eq:rd-teacher-period-spectrum}의 status에 따라 다음 경로만 허용한다.
\begin{itemize}
  \item Stable clear peak가 있으면 \(T_{\rm dom}=2\pi/\omega_{\rm dom}\)과
  \(H_s=\rho_sT_{\rm dom}\)인 period-normalized schedule을 R2b ablation 후보로 둔다.
  \item Broad/flat/multi-peak이거나 stable peak가 없으면 \(T_{\rm dom}=\mathrm{unavailable}\)이다.
  Domain-wide frozen fixed-seconds direct-long schedule을 사용하고 curriculum claim/ablation을 삭제한다.
  \item Peak가 있어도 equal-compute 이득이 없으면 curriculum contribution을 삭제하고 더 단순한 direct-long을 채택한다.
  \item 어느 경로든 final supported seconds에서 발산하면 Gate A\(_T\) prerequisite failure다.
\end{itemize}

\subsection{R2b 최소 baseline family}

R2a가 통과한 뒤 동일 data/probe/horizon/method identity에서 다음을 비교한다.
\begin{enumerate}
  \item zero/mean response,
  \item linear static-GS response regressor,
  \item capacity-matched direct set encoder without latent relation auxiliary,
  \item direct next-deformation predictor.
\end{enumerate}
Capacity/FLOPs/parameter matching rule은 test 결과 전 development config로 동결한다.

\section{구현 체크리스트}
\label{sec:r2-chapter-4}

\subsection{공통 data/model/evaluator 준비}

\begin{itemize}
  \item[$\square$] R2 slice의 object/GS/material/attachment/sequence/time grid와 R0--R1 artifact hash를 manifest에 고정한다.
  \item[$\square$] Student forward signature에 target-forbidden field가 없고 setup call count가 package당 정확히 1임을 검사한다.
  \item[$\square$] Patch-first/local evidence와 token global context, equivariant Gaussian query readout을 구현한다.
  \item[$\square$] Area head와 R0에서 동결한 predicted-or-derived validity/normal owner path를 frozen schema axis로 구현한다.
  \item[$\square$] Global field를 project/whiten한 최종 coordinate에서 pole/damping을 식별하도록 head/construction 순서를 고정한다.
  \item[$\square$] Physical-valid policy, area-derived mass, attachment projection, whitening/rank policy와 positive pole/damping을 적용한다.
  \item[$\square$] Package serialize/reload 후 같은 frame trace와 probe response를 재현한다.
  \item[$\square$] Frame-start traction, \(B_G^\top F\), one exact-ZOH update와 base mean/covariance transport만 실행한다.
\end{itemize}

\subsection{Loss/metric 구현}

\begin{itemize}
  \item[$\square$] Common-valid mask와 \(W_{A,P},W_{M,P}\)를 teacher/student 양쪽에 동일하게 적용한다.
  \item[$\square$] Displacement, velocity, FRF amplitude/phase와 rollout term마다 numerator, denominator,
  physical unit/reference scale, epsilon과 reduction axis를 \code{MetricRegistry}에 등록한다.
  \item[$\square$] Near-zero teacher amplitude에서 phase가 정의되지 않는 bin의 mask 규칙을 동결한다.
  \item[$\square$] Time sample 수, probe 수 또는 sequence 길이가 loss scale을 몰래 바꾸지 않는지 analytic duplication test를 통과한다.
  \item[$\square$] Loss weight search space, dev selection metric, final weight와 freeze 시점을 config/hash로 저장한다.
  \item[$\square$] Training loss와 acceptance metric을 구분하고 SI diagnostic과 dimensionless primary를 함께 기록한다.
\end{itemize}

\subsection{T1 --- measure/token initialization}

\begin{itemize}
  \item[$\square$] Training-only teacher quadrature/common-probe supervision으로 spatial area/measure error를 학습한다.
  \item[$\square$] Positive area와 area-derived mass, physical-valid policy 및 permutation-equivariant readout을 함께 검사한다.
  \item[$\square$] Patch-first aggregation과 global token context가 실제 gradient를 받고 nondegenerate token을 만드는지 기록한다.
  \item[$\square$] Total area/mass 합뿐 아니라 held-out training probes의 spatial measure와 analytic force/work를 fixture로 사용한다.
  \item[$\square$] T1 checkpoint/config와 T2 hand-off 시 freeze/unfreeze policy를 hash한다.
\end{itemize}

\subsection{R2a --- fixed-seconds 최소 overfit}

R2a에서는 curriculum과 baseline run을 시작하지 않는다.
\begin{enumerate}
  \item \textbf{A0 package smoke:} one-frame analytic load/update, rest identity, attachment/Gram/pole validator를 통과한다.
  \item \textbf{A1 short fixed-seconds:} impulse와 step의 \(H_{\rm smoke}\)[s] velocity/displacement/phase를 overfit한다.
  \item \textbf{A2 direct long:} 처음부터 고정한 \(H_{\rm long}\)[s]의 forced/recovery rollout과
  aero-off displaced free decay를 overfit한다.
  \item \textbf{A3 stress:} 학습보다 긴 predeclared seconds와 in-envelope force/control sweep에서 finite response와 decay ordering을 검사한다.
\end{enumerate}

\begin{itemize}
  \item[$\square$] One-step loss만 낮추고 multi-step/phase가 틀린 checkpoint를 합격시키지 않는다.
  \item[$\square$] Zero ambient rest response와 aero-off displaced free decay를 별도 initial state/flag로 평가한다.
  \item[$\square$] Force 방향 반전과 magnitude sweep에 response sign/amplitude가 일관되게 변한다.
  \item[$\square$] R1 oracle와 student를 같은 probe/mask/horizon/metric으로 비교하고 error gap을 component별로 기록한다.
  \item[$\square$] Network capacity, map/measure, response package, loss와 evaluator mismatch를 별도 failure tag로 분류한다.
\end{itemize}

\subsection{R2b --- curriculum과 baseline}

\begin{itemize}
  \item[$\square$] R2a best config/data/probe/package class를 freeze한 뒤 schedule 또는 baseline에 필요한 부분만 바꾼다.
  \item[$\square$] \(T_{\rm dom}\) available일 때 R1의 extraction site/aggregate,
  band/window/resolution/prominence와 status hash를 read-only로 검증한다.
  R2는 \(\rho_s\), transition과 final horizon만 동결해 period schedule을 실행한다.
  \item[$\square$] Period curriculum, direct-long과 fixed-frame control의 total evaluator unroll step/FLOPs를
  primary equal-compute budget으로 맞춘다. Optimizer update와 batch는 이 budget을 맞추도록 조정할 수 있다.
  \item[$\square$] Actual GPU wall time, peak memory, optimizer update, batch composition과 evaluator steps를 함께 보고한다.
  \item[$\square$] \(T_{\rm dom}\) unavailable이면 fixed-seconds fallback ID/reason을 남기고 period result를 만들지 않는다.
  \item[$\square$] Curriculum 이득이 없으면 claim과 period-specific code path를 core에서 제거하거나 optional experiment로 명시한다.
  \item[$\square$] 네 baseline을 동일 common probe, physical seconds, traction/evaluator identity와 frozen metric으로 실행한다.
  \item[$\square$] Capacity matching table에 parameter, setup FLOPs, evaluator FLOPs/steps, memory와 예외를 기록한다.
\end{itemize}

\section{Open Design Decisions}
\label{sec:r2-chapter-5}

\begin{decisionbox}{해석 규칙}
R2는 실제 loss와 optimizer를 시험하는 단계이므로 선택지가 많다. 아래 권장 기본값은 시작점일 뿐이다.
R2a acceptance run 전에 loss/metric/horizon/package-capacity를 동결하고, R2b 전에 schedule/baseline matching을
추가 동결한다. Test나 downstream Gate 결과를 본 뒤 R2 version을 소급 변경하지 않는다.
\end{decisionbox}

\begin{decisionbox}{30fps 목표의 조기 비용 확인 --- 제안, 미채택}
최종 목표는 바람 조작 시 렌더링·입력·동기화까지 약33.3 ms/frame 이내다. Teacher의 실행 시간과 별도로
R1 oracle/R2 첫 package의 힘·반응 갱신·Gaussian transport 비용을 일찍 측정하자는 의견을 제시했다.
조기 측정 범위와 시점은 아직 확정하지 않았으며, 구현 완료나 새 R2 acceptance 조건으로 취급하지 않는다.
Rank/Gaussian 수/대상 GPU와 출력 해상도를 고정하고 충분한 표현 품질을 함께 확인해야 한다.
렌더러를 제외한 비용만 측정했으면 전체30fps 달성으로 보고하지 않는다. 최종 측정은 R7 계약을 따른다.
\end{decisionbox}

\begin{longtable}{L{0.18\linewidth}L{0.26\linewidth}L{0.30\linewidth}L{0.16\linewidth}}
\toprule
결정 항목 & 허용 선택지 & 권장 기본값(미동결) & 동결 시점\\
\midrule
\endfirsthead
\toprule
결정 항목 & 허용 선택지 & 권장 기본값(미동결) & 동결 시점\\
\midrule
\endhead
R2a horizons & \(H_{\rm smoke},H_{\rm long},H_{\rm stress}\) seconds 후보 & teacher period와 무관한 domain-wide fixed seconds; short+recovery+long stress를 모두 포함 & R2a 첫 acceptance run 전\\
Global/token capacity & \(K,r_G\), patch size, attention depth/width 후보 & oracle rank와 K0 규모에 맞춘 가장 작은 nondegenerate capacity & optimizer sweep 전\\
Pole parameterization & softplus+floor / bounded log-frequency와 damping interval & domain frequency envelope를 반영한 bounded log-parameter와 positive floor & R2a config freeze 전\\
Displacement scale & raw SI / \(L_0\) / teacher per-sequence RMS & \(L_0\) 기반 domain-frozen dimensionless scale; per-sequence difficulty를 지우는 RMS는 diagnostic만 & loss search 전\\
Velocity scale & raw SI / declared \(V_{\rm ref}\) / teacher RMS & wind/time domain에서 정한 하나의 \(V_{\rm ref}\); object별 사후 scale 금지 & loss search 전\\
FRF amplitude & linear complex error / log-amplitude+phase / normalized transfer error & stabilized log-amplitude와 wrapped phase를 함께 사용하고 SI transfer diagnostic 병기 & R2a acceptance 전\\
Phase mask & fixed amplitude floor / relative SNR threshold / phase term 제외 & dev-frozen amplitude/SNR floor 아래 bin만 제외하고 mask/hash 보고 & FRF metric freeze 전\\
Probe/time reduction & sum / weighted mean / sequence-balanced nested mean & common mask의 weighted mean, sequence를 동일 outer weight로 두고 sample count duplication에 불변 & R2a acceptance 전\\
Loss weights & manual dev grid / constrained optimizer search & 좁은 log-grid와 velocity-primary dev metric; search range와 전 결과 보존 & R2a final run 전\\
Measure/relation terms & measure only / measure+relation auxiliary & exact total area/mass penalty는 빼고 spatial measure를 사용하며 relation은 diagnostic부터 시작 & R2a ablation freeze 전\\
Optimizer와 schedule & AdamW 계열 / 다른 deterministic optimizer, LR/warmup 후보 & 작은 deterministic AdamW config와 gradient/NaN trace; 성능보다 overfit 진단 우선 & R2a sweep 전\\
Period schedule & \(\rho_s\), transition, final multiple 후보 & sub-period에서 multi-period로 단조 증가, direct-long과 equal-compute & R2b run 전\\
Baseline matching & parameter / setup FLOPs / total evaluator compute 중심 & parameter와 setup compute를 함께 표기하고 rollout evaluator identity/steps는 동일 & baseline training 전\\
Acceptance threshold & oracle-relative / absolute dimensionless / 둘의 conjunction & absolute frozen metric+oracle gap diagnostic; 숫자는 dev fixture로 결정 & R2a acceptance 전\\
\bottomrule
\end{longtable}

\section{산출물}
\label{sec:r2-chapter-6}

\subsection{T1 prephase 산출물}

\begin{itemize}
  \item \code{T1InitCheckpoint}: area-measure head와 patch-to-token initialization weights,
  schema/config 및 R2a freeze/unfreeze hand-off hash.
  \item \code{T1InitReport}: spatial measure, analytic force/work, token coverage/duplicate와
  nondegenerate-gradient fixture 결과.
\end{itemize}

\subsection{R2a 산출물}

\begin{itemize}
  \item Frozen single-case dataset slice와 \code{r2a_fixed_seconds_config}, loss/metric/threshold registry revision.
  \item \code{R2aOverfitCheckpoint}: setup network checkpoint, serialized \code{ResponsePackage},
  package validation/rank/condition report와 T1 dependency hash.
  \item A0--A3 train/eval curves, common-probe trajectory/FRF, frame trace와 component failure report.
  \item Oracle--student paired comparison, fixed-seconds/free-decay/stress report와 deterministic replay command.
  \item R2a decision: pass, retry with owner-specific change, 또는 stop. Retry마다 config/version을 새로 남긴다.
\end{itemize}

\subsection{R2b 산출물}

\begin{itemize}
  \item \code{TeacherSpectrumStatus} 소비 기록과 period available/fallback branch ID.
  \item Period curriculum/direct-long/fixed-frame의 equal-compute manifest와 quality/drift/cost table.
  \item 네 baseline의 capacity/method-identity matching table, checkpoint/package와 paired result.
  \item \code{R2bHandoff}: 채택한 가장 단순한 Global training schedule/fallback branch,
  response loss, horizon, equal-compute 규칙과 frozen configuration.
  \item \code{R2GlobalCheckpoint}: R2b 종료 시 채택한 single-case Global weights/package class와
  T1/R2a/R2b dependency hash. 이는 held-out evidence가 아니라 R3 initialization artifact다.
  \item Curriculum의 유지/비주장 결정과 Gate A\(_L\)에 넘길 latent/direct-set single-case diagnostic.
  Latent-relation contribution의 최종 유지/삭제는 R3가 소유한다.
  \item \code{GateATPrerequisiteReport}: R1 teacher/map/oracle prerequisite와 R2 Global
  package/long-rollout fixture의 Boolean 결과 및 exact input hash.
  \item R2 completion report와 R3 진입 허용 여부. Single-case 결과를 held-out evidence로 표시하지 않는다.
\end{itemize}

\section{필수 fixture와 metric}
\label{sec:r2-chapter-7}

\begin{longtable}{L{0.25\linewidth}L{0.49\linewidth}L{0.16\linewidth}}
\toprule
Fixture & 합격 조건 & 단계\\
\midrule
\endfirsthead
\toprule
Fixture & 합격 조건 & 단계\\
\midrule
\endhead
Package invariant & valid-set/area/mass, attachment zero, Gram/rank, positive pole/damping과 round trip 통과 & R2a A0\\
Final-coordinate pole & projection/whitening 뒤의 field와 pole이 같은 generalized coordinate를 소유하고 raw pole 재사용이 없음 & T1/R2a A0\\
One-frame analytic & known \(F,Q\)와 exact-ZOH state/mean/covariance가 reference와 일치 & R2a A0\\
Loss duplication & 동일 probe/time sample 복제나 padding이 normalized loss를 바꾸지 않음 & R2a 전\\
Short fixed-seconds overfit & impulse/step velocity/displacement/phase가 frozen train threshold 통과 & R2a A1\\
Direct-long/recovery & \(H_{\rm long}\) forced/recovery와 aero-off free decay가 finite이고 threshold 통과 & R2a A2\\
Stress/ordering & 더 긴 seconds, force/control sweep에서 finite, sign/amplitude와 damping ordering 유지 & R2a A3\\
Oracle gap & 동일 probe/mask에서 oracle 대비 student gap과 mapping/reference floor를 component별 보고 & R2a\\
Permutation/reload & input permutation과 package reload 후 stable-ID rollout/metric이 동일 & R2a\\
Curriculum compute & available일 때 period/direct/fixed-frame의 evaluator step/FLOPs matching과 actual cost 보고 & R2b\\
Fallback identity & unavailable일 때 period result가 없고 fixed-seconds ID/reason이 manifest/table에 일치 & R2b\\
Baseline fairness & data/probe/horizon/traction/evaluator가 같고 capacity 차이가 명시됨 & R2b\\
No hidden recurrence & setup call 1회, teacher state/forcing/opaque hidden correction이 frame input에 없음 & 공통\\
\bottomrule
\end{longtable}

Primary motion metric은 velocity, displacement/tip, amplitude/phase와 target-band response다.
Renderer image quality는 이 단계의 pass condition이 아니며 base transport의 finite/rest/SPD fixture만 요구한다.

\section{종료 조건과 실패 경로}
\label{sec:r2-chapter-8}

\subsection{R2a 종료 조건}

\begin{itemize}
  \item[$\square$] Fixed loss normalization/metric/threshold와 package-capacity config가 hash됐다.
  \item[$\square$] A0 package/evaluator/transport, A1 short, A2 direct-long/free-decay와 A3 stress fixture가 통과했다.
  \item[$\square$] Attached row zero, positive pole/damping, mass Gram/rank와 deterministic replay가 통과했다.
  \item[$\square$] Oracle--student gap과 failure owner가 mapping/measure/package/loss/network/evaluator로 분리됐다.
\end{itemize}

R2a가 통과해야 R2b로 갈 수 있다. R2b는 mandatory baseline과 schedule resolution을 정리한 뒤
R3 hand-off를 발행한다. Stable peak가 없으면 period-curriculum branch만 생략하고 fixed-seconds
fallback record와 baseline은 완료한다.

\subsection{R2b 종료 조건}

\begin{itemize}
  \item[$\square$] Peak available이면 equal-compute curriculum ablation을 완료하고 유지/삭제 결정을 기록했다.
  \item[$\square$] Peak unavailable이면 fixed-seconds fallback만 실행했고 curriculum 비주장을 기록했다.
  \item[$\square$] 최소 baseline을 동일 물리/evaluator/metric 조건에서 실행하고 capacity 차이를 공개했다.
  \item[$\square$] R3가 소비할 frozen Global training config, package schema와 metric registry revision을 발행했다.
\end{itemize}

\begin{stopbox}{R2 failure와 routing}
Oracle은 통과하지만 R2a가 overfit되지 않으면 network capacity, tensor/permutation, loss normalization,
optimization과 package head를 분리 진단한다. Oracle도 같은 fixture에서 실패하면 R1의 map/measure/pole/evaluator로
돌아간다. Final supported seconds에서 NaN/divergence 또는 decay-order 붕괴가 반복되면 Gate A\(_T\) 실패다.
Dataset, Local, selector, renderer/perceptual loss를 추가해 R2 실패를 가리지 않는다.
\end{stopbox}

Curriculum이 direct-long보다 낫지 않은 것은 learned-response 가설의 실패가 아니다.
이 경우 curriculum contribution만 삭제하고 가장 단순한 fixed-seconds direct-long schedule로 진행한다.
Direct-set baseline이 latent model과 같거나 더 좋으면 R3/R4의 Gate A\(_L\) 판정에 필요한 evidence로 넘기며,
latent-anchor contribution을 자동으로 유지하지 않는다.

\end{document}
